% ============================================================================
% 6. DISCUSSION
% ============================================================================
\section{Discussion}

\paragraph{Why does Flat Text win overall?}
Flat Text provides the LLM with the \emph{complete} graph as text. For
local queries (direct FK, forward lineage), the LLM can pattern-match
the answer from the context. The 8K+ token context windows of modern
LLMs accommodate even the largest schemas. However, this advantage
vanishes on tasks requiring algorithmic traversal.

\paragraph{Why does Vector RAG underperform?}
Chunking the schema into per-table documents fragments the graph,
breaking cross-chunk edges. Despite our global lineage summary
mitigation, VR still misses schema-wide connectivity information,
explaining its poor performance on \texttt{membership} (16.7\% EM) and
\texttt{connected} (0\% EM vs.\ GA's 50\%).

\paragraph{The combined\_impact frontier.}
This subtype requires two distinct traversals: (1) forward lineage from
a source table, collecting all transitively-affected DW tables, then (2)
reverse FK from each affected table, finding all tables that depend on
them. No current baseline can reliably execute this two-phase algorithm.
We hypothesize that \emph{learned} graph representations (e.g., GNN
encoders producing per-node embeddings that capture multi-hop
reachability) are necessary.

\paragraph{Limitations.}
DW-Bench currently includes 5 datasets (4 real-world, 1 synthetic);
more diverse topologies (e.g., data mesh, lakehouse) would strengthen
generalization claims. While the Syn-Logistics synthetic dataset resolves
the statistical power limitation for low-$n$ subtypes (all subtypes
now have $n\geq20$), the real-world corpus inherently contains
small-sample subtypes. Future work should include additional models at
various scales and agentic retrieval methods.
